%==============================================================================
% Capítulo 9: Discusión
%==============================================================================
\section{Discusión}\label{sec:discusion}

\subsection{Interpretación de los resultados}\label{subsec:interpretacion-resultados}

Los resultados obtenidos confirman la hipótesis de trabajo formulada en la \autoref{subsec:hipotesis-trabajo} en lo relativo al tiempo de registro y la refutan en lo relativo a la clasificación, de modo que la hipótesis principal se declara parcialmente confirmada y parcialmente refutada: la disminución del 71,2~\% en el tiempo medio de registro es estadísticamente significativa y operativamente relevante, mientras que la exactitud estricta del 73,5~\% alcanzada por el clasificador automático quedó por debajo del umbral del 85~\% planteado en el segundo objetivo específico, resultado que se analiza entre las limitaciones de este capítulo.

La reducción de tiempo se midió con la inferencia externa incluida y con configuraciones que invocaban al modelo con mayor frecuencia que la evaluada en la clasificación, cuya etapa determinística resuelve el 65,5~\% de los casos sin invocar al modelo (\autoref{subsec:protocolo-pruebas}); esa proporción no interviene en la reducción medida. La hipótesis subsidiaria de eficiencia del esquema híbrido no se contrastó: resolver casos sin el modelo no demuestra mayor eficiencia sin comparar la exactitud de cada alternativa. En los 131 casos de esa etapa las reglas aciertan 97/131 (0,7405) y, sobre los 130 casos con predicción de ambas etapas, 96/130 (0,7385) frente a los 88/130 (0,6769) de Gemini, ejecutado con el cortocircuito deshabilitado, estimaciones ambas dentro de la muestra y sin diferencia significativa (McNemar exacta bilateral, \(p = 0{,}1516\)); además, con el vocabulario previo a su ajuste con el corpus la etapa de reglas no supera a Gemini (\autoref{subsec:sensibilidad-ajustes}), de modo que esa ventaja no se toma como evidencia de eficiencia. El costo de inferencia recae así únicamente en los casos que el filtro no resuelve, y la latencia no se midió por etapa. La hipótesis subsidiaria de integración multicanal tampoco se contrastó: no se midió la capacidad bajo carga concurrente.

La dispersión de los tiempos automatizados evidencia una compresión sustancial de la variabilidad: el desvío estándar pasa de 44,9~s en el flujo manual a 13,52~s en el automatizado, aproximadamente la tercera parte, consecuencia de la automatización y no un objetivo del diseño.

Si se proyecta el ahorro medio de tiempo por incidente (\(61{,}6 - 17{,}76 = 43{,}8\)~s) sobre el volumen trimestral de 2.646 incidentes, la automatización liberaría aproximadamente 32,2 horas-persona brutas por trimestre ---unas cuatro jornadas laborales de un operador---,\footnote{Cálculo: \(43{,}8~\text{s/incidente} \times 2.646~\text{incidentes} / 3.600~\text{s/h} \approx 32{,}2\)~h-persona por trimestre.} tiempo que el personal podría redirigir a tareas de resolución técnica o a actividades de mejora continua. Ese ahorro es bruto: no descuenta el costo de las derivaciones erróneas que el clasificador agrega respecto del flujo manual, cuyo balance se analiza en la \autoref{subsec:balance-neto}.

Sin embargo, la lectura conjunta de la intervención humana por gate y del análisis de calibración presentado en la Sección~\ref{subsec:calibracion} matiza ese resultado, y debe leerse por etapa: el gate del 70~\% presupone que la confianza declarada discrimina los casos correctos de los incorrectos, pero solo se evalúa sobre la salida de Gemini y, en el rango observado (0,85 a 1,00), esa confianza no los discrimina; la etapa determinística, por su parte, no emite una confianza utilizable para derivar. La supervisión humana efectiva no puede depender exclusivamente del umbral de confianza: debe complementarse con señales independientes del puntaje ---por ejemplo, la ambigüedad léxica entre las categorías candidatas o la revisión por muestreo de los casos clasificados con alta confianza---. El esquema \emph{human in the loop} conserva su valor de diseño como salvaguarda, pero en el corpus la regla de 0,70 no derivó ningún caso: los 124 casos con confianza registrada inferior a 0,70 pertenecen a la etapa determinística, a la que la regla no se aplica (\autoref{subsec:reduccion-intervencion}). El comportamiento del gate admite, además, una lectura propia: en los 69 casos que resolvió, la etapa Gemini nunca declaró una confianza inferior a 0,85, de modo que el gate (confianza inferior a 0,70) no derivó ningún caso; que esa conducta exprese fiabilidad o sobreconfianza no puede establecerse, porque la confianza de Gemini no discrimina aciertos de errores en el rango observado (\autoref{subsec:calibracion}).

La exactitud laxa, que acepta como acierto el sector principal predicho cuando integra el conjunto de verdad, alcanza el 89,5~\% (179/200) y se lee como la proporción de derivaciones operativamente aceptables, dado que el 53,5~\% de los casos admite más de un sector; no es la métrica del segundo objetivo específico y no modifica por sí misma su estado. Por canal, la exactitud estricta fue de 0,654 en la llamada telefónica, 0,773 en el correo electrónico y 0,811 en el formulario web (laxa: 0,864, 0,924 y 0,906), evidencia sobre la dimensión de calidad del tercer objetivo específico y el efecto potencial de la transcripción. La comparación no aísla ese efecto, porque la composición sectorial y la proporción de casos multietiqueta difieren entre canales (49 de 81 casos telefónicos frente a 15 de 53 del formulario web), por lo que la diferencia no se atribuye a la transcripción.

La procedencia de los datos de ajuste merece una observación aparte. Las reglas de frontera del prompt, la ampliación del vocabulario del filtro determinístico y el punto de operación del cortocircuito se derivaron de errores y casos observados en el corpus de evaluación; el punto de operación se fijó con un protocolo fuera de muestra (validación cruzada), pero ese protocolo no cubre la ampliación del vocabulario ni las reglas del prompt, que se definieron sobre la totalidad del corpus. No se dispuso de un corpus independiente de preproducción. En consecuencia, las métricas de clasificación reportadas son estimaciones dentro de la muestra, probablemente optimistas, cuya confirmación exige un corpus independiente (\autoref{sec:recomendaciones}). La \autoref{subsec:sensibilidad-ajustes} acota la parte de ese sesgo atribuible al vocabulario.

La reducción del tiempo de registro, aunque estadísticamente robusta, admite explicaciones rivales. Primero, la asimetría del instrumento: el flujo manual se cronometró en vivo e incluye lectura, comprensión, búsqueda y carga del ticket, mientras que el automatizado se midió con marcas de tiempo del sistema, lo que favorece al segundo; lo debilita el análisis de sensibilidad, que mantiene la reducción por encima del 65~\% aun descontando 10~s por caso, y se descartaría midiendo ambos flujos con el mismo instrumento. Segundo, la mezcla de canales: parte de la diferencia podría reflejar la composición; lo debilita el desagregado de la Tabla~\ref{tab:tiempos-canal}, con reducciones del 51,1 al 93,4~\% y contraste intra-canal significativo, aunque la magnitud global depende de la composición. Tercero, la novedad o efecto Hawthorne y el aprendizaje del operador: los operadores sabían que se cronometraba y el procedimiento lo atenúa sin eliminarlo (\autoref{subsec:diseno-experimental}); el orden temporal de la campaña (manual antes que automatizado) sesga, de hacerlo, en contra de la hipótesis, y se descartaría con un diseño contrabalanceado. Cuarto, la selección de los 200 casos: por ser una muestra estratificada de una sola organización, la magnitud observada no es extrapolable sin nuevas validaciones. Ninguna de estas explicaciones rivales se descarta por completo con el diseño actual.

El patrón de error más informativo es la confusión sistemática de Sistemas: de los 16 casos asignados, solo 7 se clasificaron correctamente, y los errores se repartieron entre Soporte Tecnico Software (5), Bases de Datos (2) y Seguridad Informatica (2). Esta frontera combina dos causas que conviene no confundir. La primera es una ambigüedad léxica y taxonómica: las descripciones de acciones sobre servidores, servicios de plataforma o aplicaciones admiten más de un sector, y la etiqueta de referencia resuelve esa ambigüedad mediante una convención que el clasificador no conoce. La segunda es falta de contexto: muchas entradas son demasiado breves para desambiguar el sector responsable, y el modelo carece de señales que un operador obtendría por conocimiento organizacional. Solo la segunda admite mitigación por enriquecimiento del protocolo de captura; la primera exige revisar la taxonomía o el criterio de etiquetado. Cabe advertir, además, que Bases de Datos (\(n = 1\)) no es estimable: su F1 no debe leerse como evidencia de que el clasificador falla en esa clase, sino como ausencia de soporte estadístico.

Esa lectura no equivale a una confirmación concluyente: la evidencia es compatible con explicaciones rivales que el diseño no permite descartar por completo. Mientras que la reducción observada admite una lectura causal fuerte en el plano práctico ---el sistema registra cada incidente en menos tiempo que el proceso manual---, su atribución a cada componente del esquema y su extrapolación a otras organizaciones permanecen abiertas. Las amenazas a la validez que acotan estas conclusiones están enumeradas en la \autoref{subsec:amenazas-validez}.

\subsection{Comparación con la literatura}\label{subsec:comparacion-literatura}

Los valores de exactitud y F1 macro obtenidos se contrastan a continuación con los reportados en la literatura comparable. \textcite{Paramesh2019_it_service_desk} comparan Naive Bayes multinomial, regresión logística, k vecinos más cercanos y SVM sobre tickets reales y reportan buen desempeño de SVM en todas las muestras, sin que el resumen consigne cifras comparables. \textcite{Revina2020_it_tickets}, en cambio, muestran que los clasificadores simples y explicables construidos con rasgos lingüísticos del dominio son competitivos frente a alternativas más complejas.

Los resultados del presente trabajo ---73,5~\% de exactitud estricta y un F1 macro de 0,6508 sobre las cuatro clases estimables (0,5207 al incluir la clase de soporte unitario, no estimable)--- no admiten una comparación directa con la literatura revisada: los trabajos citados reportan clasificación de etiqueta única, mientras que la evaluación propia se mide contra una referencia multietiqueta sobre cinco sectores canónicos en español rioplatense, tarea considerablemente más exigente. La exactitud laxa (89,5~\%) no es comparable con las cifras de etiqueta única de aquellos trabajos y se informa solo como referencia secundaria. Frente al enfoque SVM de \textcite{Paramesh2019_it_service_desk} y a la evidencia de \textcite{Revina2020_it_tickets} sobre la competitividad de los métodos simples, el esquema propuesto no requiere infraestructura de GPU para inferencia local ni un corpus de entrenamiento etiquetado de gran escala, lo que lo hace más accesible para organizaciones medianas, aunque la evidencia disponible no permite ubicar su exactitud respecto de la de esos enfoques.

La evidencia de \textcite{Revina2020_it_tickets} resulta, en apariencia, contraria al enfoque del presente trabajo, que apuesta por un LLM de propósito general. Esta tensión, sin embargo, refuerza el análisis: aquellos autores excluyen por diseño las representaciones preentrenadas, al considerar que el vocabulario específico del dominio podría degradarlas, lo que desaconseja suponer una transferencia no supervisada trivial. Precisamente por eso la elección del LLM en este trabajo no se justifica por transferencia pasiva de conocimiento lingüístico, sino por la tarea concreta ---clasificación de categoría con contexto instruccional y esquema híbrido de reglas más modelo---, en la que el modelo se emplea como motor de inferencia guiado y no como representación preentrenada sin ajuste.

La literatura de calibración y cascadas (\autoref{subsec:calibracion-cascadas}) permite leer dos resultados. La confianza de la etapa Gemini se concentró en el extremo superior de la escala y no discriminó aciertos de errores en su rango observado (\autoref{subsec:calibracion}), en línea con la sobreconfianza y la baja discriminación que \textcite{Xiong2024_confidence_elicitation} documentan para la confianza verbalizada, aunque \textcite{Tian2023_verbalized_confidence} muestran que depende de la elicitación. En las cascadas la derivación exige una señal de fiabilidad correlacionada con el acierto y un umbral \parencite{Varshney2022_model_cascading, Chen2024_frugalgpt}: el cortocircuito de la etapa determinística la adopta con el puntaje de correctitud, que ordena la correctitud esperada sin ser una probabilidad calibrada, mientras que la revisión humana depende de la confianza verbalizada de Gemini, la señal que esa literatura describe como sobreconfiada. La eficiencia del esquema híbrido frente a una alternativa solo-modelo sigue sin contrastarse (\autoref{subsec:interpretacion-resultados}).

La novedad del aporte radica en la validación específica sobre español rioplatense aplicado a un dominio de mesa de ayuda de organización mediana, segmento escasamente cubierto por la literatura previa. Diversos trabajos han señalado que el ajuste y la evaluación de los LLM se concentran en el inglés y que su rendimiento cae en los idiomas con menos recursos \parencite{lai2023okapi, xuan2025mmluprox}, y los resultados del presente trabajo (F1 macro = 0,6508 sobre las cuatro clases estimables) sugieren que el desempeño en español rioplatense resulta inferior al reportado para el inglés, brecha que se atribuye, al menos en parte, al esquema multietiqueta de cinco sectores adoptado en esta evaluación.

\subsection{Limitaciones del estudio}\label{subsec:limitaciones-estudio}

El estudio presenta ocho limitaciones que deben considerarse al interpretar los hallazgos y al planificar replicaciones futuras.

Primero, el tamaño de la muestra (\(n = 200\)) es suficiente para una validación piloto pero insuficiente para sustentar generalizaciones de carácter poblacional sobre el universo de mesas de ayuda en organizaciones medianas. El intervalo de confianza al 95~\% reportado para la exactitud estricta, [67,0~\%, 79,1~\%], se ubica íntegramente por debajo del umbral del 85~\%, por lo que el segundo objetivo específico no se alcanza; muestras mayores, un corpus más equilibrado entre clases y una revisión del clasificador serían necesarios para cerrar esa brecha.

Segundo, la circunscripción del corpus a una única organización del sector servicios de la provincia de Mendoza restringe la validez externa de los hallazgos. La distribución de tipos de incidentes, el léxico organizacional y los patrones de redacción pueden variar significativamente entre organizaciones, sectores y regiones.

Tercero, la dependencia operativa de un servicio externo de inferencia condiciona el rendimiento del sistema. Un eventual deterioro del servicio de Gemini~3.6 Flash, un cambio significativo en las condiciones comerciales del proveedor, o una modificación en el comportamiento del modelo debido a actualizaciones afectaría directamente la operación. Esta dependencia es intrínseca al enfoque basado en LLM externo y debe ser gestionada como un riesgo operativo.

Cuarto, los resultados reflejan el comportamiento del modelo en su versión vigente durante la campaña automatizada (\texttt{gemini-3.6-flash}, integrada al clasificador el 24 de septiembre de 2026; \autoref{fig:cronograma}). Las actualizaciones futuras del proveedor pueden modificar tanto el desempeño bruto como la estructura del prompt requerido para alcanzar resultados equivalentes, lo que aconseja la implementación de un régimen de monitoreo continuo de las métricas de clasificación.

Quinto, la comparación se realizó exclusivamente contra el flujo manual (proceso operativo), no contra clasificadores alternativos como TF-IDF+SVM o regresión logística sobre el mismo corpus. Un chequeo exploratorio temprano con un clasificador TF-IDF+SVM sobre un subconjunto no preservado del corpus no se reporta aquí como comparación válida: ni la partición ni el tamaño del subconjunto quedaron registrados, por lo que su cifra no es reproducible ni comparable con las métricas de este capítulo. Se lo menciona únicamente como antecedente del diseño, no como evidencia de superioridad del esquema híbrido. Una comparación sistemática contra clasificadores supervisados sobre el corpus completo, con partición y código preservados, queda como trabajo futuro.

Sexto, los analistas que etiquetaron la referencia operaron también el brazo manual (\autoref{subsec:poblacion-muestra-corpus}), lo que acota la independencia del \emph{ground truth} y, junto con la falta de registro por caso de la derivación manual, impide medir sin sesgo la exactitud de ese brazo (\autoref{subsec:balance-neto}).

Séptimo, el vocabulario del filtro determinístico y las reglas de frontera del prompt se ajustaron con errores observados en el mismo corpus de evaluación, de modo que las métricas de clasificación son estimaciones dentro de la muestra. Sin ese vocabulario, la exactitud de la etapa determinística baja de 0,7405 a 0,5270, cota pesimista del sesgo atribuible a él, y el efecto de las reglas del prompt, acotado por una ablación controlada de dos brazos, es de 50 frente a 46 aciertos sobre los 69 casos de la etapa Gemini (diferencia no significativa; McNemar exacto \(p = 0{,}3438\); \texttt{evaluation/report\_sensibilidad\_ajustes.md}, línea~98) (\autoref{subsec:sensibilidad-ajustes}). Solo un corpus independiente permitiría estimar el desempeño fuera de la muestra.

Octavo, no se conservó un registro primario de la acción manual de los operadores por caso, de modo que la proporción de intervención humana no se midió y la segunda sub-hipótesis queda no contrastada (\autoref{subsec:reduccion-intervencion}).

\subsection{Implicancias prácticas}\label{subsec:implicancias-practicas}

Desde el punto de vista práctico, los resultados sugieren que organizaciones medianas con volúmenes de incidentes comparables al observado en la organización analizada pueden reducir sustancialmente el tiempo de registro mediante soluciones híbridas similares a la propuesta, con un saldo neto que depende de la exactitud efectiva (\autoref{subsec:balance-neto}). La automatización de la etapa de registro puede liberar al personal técnico para concentrarse en la resolución efectiva de los casos, etapa en la cual el juicio experto continúa siendo insustituible.

La trazabilidad provista por el registro detallado de cada decisión del clasificador habilita dos procesos de mejora continua. Por un lado, permite identificar patrones de error sistemáticos y ajustar el filtro determinístico o el prompt del LLM en consecuencia. Por otro lado, la acumulación sostenida de decisiones validadas por los sectores responsables constituye un activo informacional que, una vez alcanzado un volumen adecuado del orden de 5.000 casos (referencia indicativa de orden de magnitud, no derivada del corpus actual ni de un cálculo de potencia estadística), permitiría entrenar modelos de clasificación específicos del dominio organizacional, con eventuales ventajas en costo, latencia y soberanía sobre los datos.

En el plano económico, el costo operativo del sistema se compone de dos rubros principales. El primero es la infraestructura autoalojada: una instancia pública con 2 vCPU y 4 GB de memoria se ofrece a 24 USD/mes, rubro que incluye 4.000 GiB de transferencia saliente y 80 GiB de disco SSD y al que se suman, si se contratan, el respaldo automático (entre el 20~\% semanal y el 30~\% diario del costo de la instancia) y otros servicios de almacenamiento y direccionamiento; el costo mensual se ubica así en el orden de 24 a 32 USD/mes \parencite{DigitalOcean2026_pricing}.\footnote{Precio de lista consultado el 19 de octubre de 2026 en \url{https://www.digitalocean.com/pricing/droplets}: Droplet Basic de 2 vCPU y 4 GB, 24 USD/mes, con 4.000 GiB de transferencia y 80 GiB de SSD incluidos; el respaldo automático agrega entre el 20~\% (semanal) y el 30~\% (diario) del costo de la instancia. No se computan volúmenes adicionales, instantáneas ni direcciones IP reservadas.} El segundo rubro es el costo de inferencia del LLM. Con el precio introductorio vigente del modelo \texttt{gemini-3.6-flash} (0,75 USD por millón de tokens de entrada y 3,75 USD por millón de tokens de salida) \parencite{Google2026_vertex_pricing},\footnote{Precios de lista consultados el 19 de octubre de 2026 en \url{https://cloud.google.com/vertex-ai/generative-ai/pricing}: Gemini 3.6 Flash, precio introductorio de 0,75/3,75 USD por millón de tokens de entrada/salida vigente hasta el 31 de diciembre de 2026; desde el 1 de enero de 2027 rige el precio estándar de 1,50/7,50 USD por millón de tokens, que duplica la cifra. Equivalen a \(0{,}75\times10^{-6}\) y \(3{,}75\times10^{-6}\) USD por token.} y con supuestos explícitos de aproximadamente 800 tokens de entrada (prompt con la descripción de las cinco categorías, ejemplos balanceados y el incidente) y 30 tokens de salida (JSON estricto con categoría y confianza), cada invocación cuesta del orden de \(7\times10^{-4}\) USD, entre 0,0006 y 0,0009 USD según el tamaño del prompt.\footnote{Cálculo: \(800\times10^{-6}\times0{,}75 + 30\times10^{-6}\times3{,}75 \approx 7\times10^{-4}\) USD por invocación; el límite de salida configurado es de 100 tokens y el JSON de respuesta ronda los 15--30 tokens. No se computan reintentos ni la transcripción de voz, que se factura por separado.} Para un volumen de 2.646 incidentes trimestrales, con un 34,5~\% requiriendo inferencia externa, el costo trimestral de API se sitúa en el orden de 1 USD y en ningún caso por encima de 2 USD, un valor marginal frente a las 32,2 horas-persona brutas liberadas.

Un aspecto que excede el alcance técnico de la solución pero condiciona su efectividad en operación real es la dimensión organizacional de la adopción. La experiencia de campo indica que los empleados con mayor antigüedad tienden a resistir los cambios de procedimiento, incluso cuando ofrecen beneficios demostrables, por familiaridad con los procedimientos existentes y percepción ---fundada o no--- de que el nuevo sistema puede resultar complejo o invasivo. Por eso la puesta en producción se beneficia de un despliegue gradual: un grupo piloto reducido, extensión progresiva a otros sectores y capacitación breve y contextualizada en cada etapa. El enfoque incremental reduce la fricción inicial, permite ajustar la configuración a cada sector y facilita la confianza de los usuarios finales, condición necesaria para que los beneficios cuantificados ---reducción del 71,2~\% en el tiempo de registro--- se traduzcan en mejoras organizacionales sostenibles.

\subsection{Balance neto de la automatización}\label{subsec:balance-neto}

La proyección de 32,2 horas-persona por trimestre es un ahorro bruto: no descuenta las derivaciones erróneas que el clasificador agrega respecto del flujo manual. Sean \(N = 2.646\) incidentes por trimestre (\autoref{subsec:interpretacion-resultados}), \(\Delta t = 43{,}82\)~s el ahorro medio por incidente (\autoref{sec:resultados}), \(e_{\mathrm{aut}} = 0{,}265\) (53/200) la tasa de error estricta del clasificador, \(e_{\mathrm{man}}\) la tasa de derivación errónea manual (15~\% preliminar; \autoref{subsec:planteamiento-problema}) y \(c\) el costo de cada una (10 a 15 minutos, estimación propia; \autoref{subsec:objetivos-especificos}). El balance neto trimestral, en horas-persona, es \(B = (N/3.600)\,[\Delta t - (e_{\mathrm{aut}} - e_{\mathrm{man}})\,c]\) y la exactitud mínima de equilibrio, \(a^{*} = 1 - e_{\mathrm{man}} - \Delta t/c\), independiente de \(N\).

\begin{table}[htbp]
\centering
\caption{Exactitud mínima de equilibrio y balance neto trimestral según la tasa de derivación errónea manual y el costo por derivación}
\label{tab:balance-neto}
\footnotesize
\begin{tabular}{@{}l ccc ccc@{}}
\toprule
 & \multicolumn{3}{c}{\textbf{Exactitud mínima \(a^{*}\) (\%)}} & \multicolumn{3}{c}{\textbf{Balance neto \(B\) (h-persona)}} \\
\cmidrule(lr){2-4}\cmidrule(lr){5-7}
\textbf{Tasa manual} & 5 min & 10 min & 15 min & 5 min & 10 min & 15 min \\
\midrule
10~\% & 75,4 & 82,7 & 85,1 & \(-\)4,2 & \(-\)40,6 & \(-\)76,9 \\
15~\% (preliminar) & 70,4 & 77,7 & 80,1 & +6,9 & \(-\)18,5 & \(-\)43,9 \\
20~\% & 65,4 & 72,7 & 75,1 & +17,9 & +3,5 & \(-\)10,8 \\
26,5~\% (paridad) & 58,9 & 66,2 & 68,6 & +32,2 & +32,2 & +32,2 \\
\bottomrule
\end{tabular}
\par\smallskip
{\footnotesize\textit{Fuente:} elaboración propia a partir de \(N = 2.646\), \(\Delta t = 43{,}82\)~s y \(e_{\mathrm{aut}} = 53/200\). Tasas y costos son supuestos de sensibilidad, no mediciones. Referencias: exactitud estricta 73,5~\% (IC 95~\%: [67,0~\%; 79,1~\%]) y laxa 89,5~\% (IC 95~\%: [84,5~\%; 93,0~\%]). Una decimal; con \(\Delta t = 43{,}8\)~s dos celdas difieren en 0,1.\par}
\end{table}

Con 15~\% y entre 10 y 15 minutos por reasignación, el clasificador agrega unas 304 derivaciones erróneas por trimestre, que cuestan entre 50,7 y 76,1 horas-persona frente a las 32,2 de ahorro: un balance de \(-\)18,5 a \(-\)43,9 horas-persona. El equilibrio tolera entre 7,3 y 4,9 puntos porcentuales de error adicional; el observado es de 11,5. La \autoref{tab:balance-neto} muestra que el signo depende de los supuestos. El intervalo de Wilson de la exactitud estricta, [67,0~\%; 79,1~\%], contiene a \(a^{*} = 77{,}7\)~\% (15~\%, 10 minutos), de modo que ni con esos supuestos la muestra fija el signo; con 15 minutos, \(a^{*} = 80{,}1\)~\% lo supera. El umbral del 85~\% equivale a la paridad con el 15~\% manual, y \(a^{*}\), que solo acredita el ahorro de registro, es menos exigente. Con la exactitud laxa, que supera los doce valores de \(a^{*}\), el balance resultaría positivo en todos los escenarios (de +28,9 a +138,0 horas-persona; +42,1, +52,1 y +62,0 en el central); es, sin embargo, la lectura favorable, porque supone que ninguna derivación a un sector válido genera reasignación.

El balance neto es, entonces, indeterminado con el diseño actual. La exactitud del brazo manual no se midió: el corpus no registra por caso la derivación del operador y los operadores son los mismos analistas que fijaron la referencia, de modo que una medición sobre este corpus estaría sesgada hacia el acuerdo. El 15~\% es una estimación preliminar de la misma naturaleza que los 2 minutos 45 segundos luego corregidos a 61,6~s, y el 26,5~\% es una cota superior de la tasa operativa, porque el 53,5~\% de los casos admite más de un sector. Quedan así sostenidas la viabilidad técnica (reducción del 71,2~\% del tiempo de registro) y la económica (costo monetario marginal), mientras que la operativa es condicional a que la exactitud efectiva supere \(a^{*}\), por mejora del clasificador o porque la revisión humana del sector responsable reduzca la tasa o el costo de las derivaciones erróneas. Medir ambos es una línea futura (\autoref{subsec:estudios-replicacion}).

\subsection{Reflexiones sobre la generalización}\label{subsec:reflexiones-generalizacion}

La generalización de los resultados a otras organizaciones requiere atender al menos tres condiciones. En primer lugar, la disponibilidad de un corpus etiquetado de tamaño y calidad similares al utilizado en este trabajo, lo cual implica una inversión inicial en doble etiquetado y validación del acuerdo inter-evaluador. En segundo lugar, la viabilidad técnica de mantener desplegada localmente la infraestructura de orquestación y persistencia, condición que descarta a organizaciones sin capacidades operativas mínimas para la administración de contenedores Docker. En tercer lugar, la existencia de un marco normativo compatible con el envío de fragmentos descriptivos de incidentes a un servicio externo de inferencia, marco que en el caso argentino se encuentra contemplado por la Ley~25.326 \parencite{Ley25326_2000} bajo determinadas condiciones de minimización y pseudonimización descritas en la \autoref{sec:aspectos-legales}.

\newpage
